% Generated by scripts/build_textbook.py; edit content/ and rebuild.

\chapter{High-Velocity Delivery}

For decades, arguments about software delivery were settled by whoever had the most seniority in the room. Ops wanted fewer releases because releases caused outages; developers wanted more releases because customers were waiting. Both sides had anecdotes. Neither had evidence.


\index{DORA metrics}
That changed when the DORA research programme examined software-delivery practices and outcomes across industries. Its original four-key model became a common language for delivery performance; the current model uses five measures and distinguishes change throughput from deployment instability. The durable finding is that speed and stability are not general trade-offs: top performers tend to do well across the measures because small changes, automated pipelines and quick recovery support both. \autocite{dora-research,dora-metrics}


\index{CI/CD}
\index{error budget}
\index{GitHub Actions}
This part is about those habits. In the book's through-line, this is where the service starts changing safely and repeatedly rather than depending on heroic manual releases. You'll learn the current five DORA metrics and how to interpret them, how a CI/CD pipeline moves a commit safely to production, how to build one in GitHub Actions, why branch lifetime matters more than branch strategy, and how error budgets let engineers and product managers argue about risk with numbers instead of volume. \autocite{github-workflows,google-error-budget}


\index{DevOps}
\index{platform engineering}
\index{SRE}
\index{YAML}
Day to day, DevOps and SRE work feels less glamorous than the job ads suggest: editing YAML, watching a pipeline go red, reading logs, keeping merges small, occasionally being woken by a pager. But it's also where a junior engineer can have outsized impact, because every hour of toil you automate away is returned to the whole team, every week. The part closes with a look at the DevOps, SRE and platform engineering career paths — some of the best-paid, most remote-friendly roles in the industry.


In the hands-on component you'll build a small pipeline, deliberately break it, and measure your own lead time and recovery time. By the end, high delivery performance should feel less like a slogan and more like a set of visible habits.

\section{CI/CD Pipeline Design}
\index{CI/CD}
In Part 6 you'll meet Sarah, whose fifteen-person marketplace startup ships to customers every week. Early on, ``shipping'' meant one of her two engineers spending Friday afternoon copying files onto the production server by hand, running a half-remembered sequence of commands, and hoping. Twice the sequence was run out of order. Once a config file was skipped entirely, and the site served a broken checkout page all weekend because nobody thought to look. The problem wasn't carelessness — it was asking humans to be perfectly repeatable, which is the one thing humans are reliably bad at.


Continuous integration and continuous delivery (CI/CD) exist to end exactly that pain. \textbf{Continuous integration} means every code change is automatically built and tested the moment it lands, so problems surface in minutes rather than at release time. \textbf{Continuous delivery} means the software is kept in a releasable state, with deployment itself automated down to a button-press (or no press at all). Together they form a pipeline: an automated assembly line that carries a change from a developer's commit to running in production, applying the same checks in the same order every single time.



\subsection{Why pipelines exist}


\index{DORA metrics}
Manual deployments fail for predictable reasons. They're slow, so teams batch changes up and release rarely — which, as the DORA research showed, makes each release riskier. They depend on one person's memory, so knowledge concentrates in whoever ``does the deploys'' and evaporates when that person is on leave. And they're invisible: when the release breaks, reconstructing what actually happened means archaeology through shell histories.


Automation reverses each of these. A pipeline makes releases repeatable — the same steps, in the same order, with a log of every run. It makes them fast and cheap, so deploying twice a day costs no more courage than deploying twice a year. It gives quick feedback: a developer who broke the build finds out in ten minutes, while the change is still fresh in their head, instead of three weeks later during integration testing. And it returns engineers' attention to the product. Sarah's startup doesn't employ anyone whose job is ``shipping chores'' — the pipeline does that, and it doesn't take Fridays off.


This is not just a big-company practice. Small teams arguably benefit most, because they have no spare capacity to absorb a lost weekend. Even non-technical staff feel the difference: instead of asking whether the fix ``actually made it to the customer site,'' anyone can look at the pipeline status and know.



\subsection{Four design principles}


Well-designed pipelines vary enormously in tooling but share a small set of principles.


\begin{itemize}[leftmargin=*]
\item \textbf{Automate the whole path.} Build, test and deploy steps all run without manual intervention. Anything a human must remember to do will eventually be forgotten.
\item \textbf{Build once, deploy many times.} Compile or package the application exactly once, version that artifact, and promote \emph{the same artifact} through test, staging and production. If you rebuild for each environment, you are no longer testing what you ship.
\item \textbf{Fast feedback first.} Order the pipeline so the quickest checks run earliest. A developer should hear about a failing unit test in minutes, not after an hour-long deployment rehearsal.
\item \textbf{Gates, not gatekeepers.} Security scans and quality checks are wired into the pipeline itself, so nothing ships without passing them — and nobody has to be the bad guy who blocks a release, because the pipeline does it impartially.
\end{itemize}

\index{containers}
The second principle deserves emphasis because it's the one beginners violate most. The artifact (a container image, a zip of compiled code, a versioned package) is the unit of trust. Once it has passed testing, promoting that exact artifact means production runs the thing you validated, byte for byte.



\subsection{From commit to production}


A typical pipeline runs as a sequence of stages, each of which either builds confidence in the change or stops it early.


It begins when a developer pushes a \textbf{commit} (code and configuration together) to source control. That event triggers the pipeline automatically; nobody has to remember to start it. The \textbf{build} stage spins up a clean, disposable environment and produces the versioned artifact there. The clean environment matters: leftover files from a previous build are a classic source of ``works in CI, fails in production'' mysteries. The \textbf{test} stage runs the automated suite (unit tests, integration tests, linting) against the artifact. A \textbf{security} stage scans dependencies for known vulnerabilities and enforces policy gates, so basic health checks can never be skipped in the rush to release.


Only then does \textbf{deployment} begin, and it proceeds progressively rather than in one leap. The artifact goes first to a staging environment that mirrors production as closely as possible; that's where smoke tests run and where a human approval can be required if the team wants one. When staging looks good, the same artifact is promoted to production. Finally (and this stage is skipped surprisingly often) the pipeline's job continues after release: an \textbf{observe} stage watches metrics, logs and alerts for signs that production is degrading. If the signals turn bad, the escape path is simple and fast: roll back to the previous release, which is still sitting there as a known-good artifact. A pipeline without a rehearsed rollback path is a one-way door.



\begin{quote}\small
The whole flow is worth memorising as a chant: commit, build, test, security, deploy, observe — with a return arrow from observe back to the previous release. Every stage either builds confidence or stops the change early. That's the entire design philosophy in one sentence.
\end{quote}



\subsection{What this looks like in practice}


\index{GitHub Actions}
\index{YAML}
The most accessible way to build one of these today is GitHub Actions, which the next topic covers in detail. The short version: you describe the pipeline in a YAML file that lives in the repository itself, triggered whenever code is pushed. Separate jobs handle building, testing and deploying, which keeps failures easy to localise — when the run goes red, you can see at a glance which stage failed. Common tasks (checking out code, setting up a language runtime, uploading artifacts) use reusable actions shared by the community, and credentials live in encrypted secrets rather than being pasted into scripts. The pipeline gets a locked safe for its keys; your git history stays free of passwords.



\subsection{The takeaway}


A pipeline is less a piece of tooling than a trust-building machine. Every green run is a small proof that the path from laptop to production works; over hundreds of runs, that proof compounds into a team that deploys on a Tuesday afternoon without holding its breath. Mistakes still happen — pipelines catch many of them before customers notice, and make the rest cheap to reverse.


The practical advice is to invest early. A pipeline built in week one of a project costs an afternoon; retrofitting one onto two years of manual-deployment habits costs a quarter. Sarah's startup got there after the broken-checkout weekend, and the pattern held: faster feedback, fewer surprises, and calmer releases — moved from midnight to mid-morning, because there was no longer any reason to hide them in the quiet hours.

\section{DevOps, SRE \& Platform Careers}
\index{DevOps}
\index{SRE}
\index{DevOps engineer}
\index{Kubernetes}
\index{platform engineering}
\index{SRE engineer}
Open Seek on any weekday and search for ``Kubernetes''. Among the results you'll find three job titles that look, on first read, like the same job wearing different hats: DevOps engineer, site reliability engineer, platform engineer. All three ads mention pipelines, cloud, automation and an on-call roster. All three pay well. And all three barely existed fifteen years ago — tell someone in 2010 that you were a DevOps engineer and there was a fair chance they'd ask if you fixed printers. The titles condensed out of the same pressure: companies like Google and Netflix demonstrated that at serious scale you cannot ship quickly \emph{and} stay reliable without automation and tight feedback loops, and whole professions formed around that discovery.


The roles overlap heavily (same tools, adjacent desks, often the same incident channel) but each leans in a distinct direction. Working out which lean suits you matters more than memorising a title. Salary bands move with country, city, remote policy, seniority and specialisation, so verify current local advertisements and salary surveys rather than relying on a printed figure. The steadier pattern is that these roles are valued where outages, manual toil and developer waiting time are expensive.



\subsection{DevOps engineers: the bridge builders}


\index{GitHub Actions}
A DevOps engineer's territory is the pipeline — everything between a developer typing \texttt{git push} and the change running in production. A normal day might involve wiring up GitHub Actions or Jenkins jobs, containerising an application with Docker, and occasionally debugging a failed deployment at an hour when reasonable people are asleep. The unofficial motto of the trade: ``it works on my machine'' is not a release strategy.


\index{DORA metrics}
\index{YAML}
The defining skill isn't any single tool, though; it's the bridging. DevOps engineers sit between development and operations and translate. Developers want speed, operations wants stability, and (as the DORA research covered in this part keeps showing) a good pipeline is how you deliver both at once. That makes communication as load-bearing as scripting. The people who thrive in the role are curious, comfortable with ambiguity, and able to explain a build failure to someone who has never seen a line of YAML.


\index{AWS}
\index{open source}
Almost nobody studies for this job directly. Most arrive from one of two directions: system administrators who discover a knack for automation, or developers who keep volunteering to fix the deploy scripts until it becomes their whole job. Certifications such as AWS Certified DevOps Engineer – Professional or the Kubernetes CKAD can be useful progress markers, but names and employer preferences change, so verify the current requirements. Remote work and on-call rotations are both common patterns, though neither should be assumed from a title. A small startup might have a single DevOps generalist covering everything from CI to cloud billing; a large enterprise will run teams with room to specialise. From here the paths lead to architecture, engineering management or deep involvement in the open source tooling the industry runs on.



\subsection{Site reliability engineers: guardians of the numbers}


\index{error budget}
\index{SLO}
Site reliability engineering is a Google invention, born from the problem of running systems too large for humans to babysit. The founding idea (treat operations as a software problem) produced the error budgets and service level objectives you met earlier in this part, and SREs are the people who hold teams to them.


The day-to-day mixes engineering and vigilance. An SRE might spend the morning writing a Kubernetes operator, the afternoon tuning Prometheus alert thresholds, and the week after that running a chaos experiment — deliberately killing components to prove the system survives, a practice Netflix pioneered after its own early outages. When an incident does land, the SRE is often the calm voice running the response, and afterwards the person leading the retrospective. Two temperamental traits matter more than any qualification: staying level when the 3 am page arrives, and genuinely liking metrics. SRE is the most numbers-driven role in this trio; if the phrase ``our error budget is 43\% consumed'' excites rather than bores you, that's a signal.


Most SREs come from software development or DevOps backgrounds and are comfortable both writing code and debugging production under pressure. Cloud and reliability certifications can help as credibility markers, but current employer preferences change quickly. Reliability specialists are expensive because outages are more expensive. On-call rotations and incident retrospectives are simply part of the culture, though many teams are remote-friendly. Large companies maintain dedicated SRE squads; smaller firms rely on one or two specialists who watch the dashboards closely because the business depends on them.



\subsection{Platform engineers: paving the roads}


If DevOps engineers build the pipeline for a team and SREs keep the service standing, platform engineers build the roads everyone drives on. Their customers are internal: the other developers in the company. The output is ``paved roads'' — reusable infrastructure modules, standard Kubernetes templates, golden machine images, self-service environments a developer can spin up without filing a ticket, and internal developer portals. Spotify's Backstage is the canonical example: an internal portal that worked so well it was open sourced and is now used across the industry.


Platform engineers usually arrive from DevOps or SRE roles, having noticed they enjoy building shared services more than nursing a single application. The extra ingredients are empathy and product thinking — a platform nobody wants to use is just infrastructure with better marketing, so good platform engineers interview their internal users, watch where developers get stuck, and prioritise like product managers. Large enterprises may run platform teams of twenty or more; a startup might have one person quietly keeping every team productive. The career path bends towards platform architect or product management roles, which makes this the trio's most natural bridge out of pure engineering.


It is also the least visible job of the three when done well — nobody notices the road builders until there's a pothole. If you need regular public credit for your work, that's worth knowing about yourself before you take the role.



\subsection{Choosing a lane — and changing it}


Company size shapes what these jobs actually look like. In a fifteen-person startup the titles blur into one generalist who does all three badly on Tuesday and brilliantly on Wednesday. Mid-sized companies carve out dedicated SRE or platform functions once the pain justifies it. Large enterprises can support entire divisions devoted to reliability or internal tooling. None of these is the ``real'' version; they're the same skills at different scales.


A rough temperament test: if you love automation \emph{and} translation (smoothing the friction between teams), lean DevOps. If you love measurement and are unusually calm in a crisis, lean SRE. If you love building tools whose users sit two desks away, lean platform. The boundaries are porous, the skills transfer almost completely, and plenty of careers touch all three. The shared path runs from hands-on engineer to strategic leader. Pay, remote-work policy and on-call compensation vary substantially by market and seniority; use current local advertisements and salary surveys when comparing offers. Certifications in AWS, GCP or Kubernetes can help at the junior end; at the senior end, a history of reliable systems speaks louder.



\begin{quote}\small
When you're the one being interviewed, ask them: ``What does your on-call rotation actually look like — frequency, compensation, and how often does it page?'' The answer tells you more about the job than the title does, and how comfortably they answer tells you about the culture.
\end{quote}


These fields also have a diversity problem. Participation from underrepresented groups remains uneven, and broad claims about improvement should be checked against current workforce data. Whether you specialise in one lane or blend all three, employers still need people who can automate delivery, reason about reliability and explain the machinery to others.

\section{DORA Metrics}
\index{DORA metrics}
Picture the monthly release meeting at Meridian Insurance, a mid-sized firm with a claims platform that most of its two thousand staff touch every day. The head of operations wants to push the next release back a fortnight because the last one caused a Saturday outage. The development manager wants to ship this week because three teams have finished features promised to customers. Both are certain they are right, and the meeting is drifting towards a decision by seniority and stamina.


For much of the industry's history, that was the state of the argument. Speed and stability were treated as opposite ends of a rope, and organisations chose a position through instinct and recent trauma. The DORA research programme matters because it gave teams a way to test that story with delivery data. DORA's research and definitions continue to evolve, so use its current guidance rather than memorising a diagram from an old report. \autocite{dora-research,dora-metrics}



\subsection{From The Original Four Keys To Five Metrics}


\index{change failure rate}
\index{DevOps}
\index{DORA metrics!change failure rate}
\index{DORA metrics!deployment frequency}
\index{DORA metrics!lead time for changes}
DORA began as DevOps Research and Assessment, associated with researchers including Nicole Forsgren, Jez Humble and Gene Kim. Its long-running surveys examined software-delivery practices and organisational outcomes across industries. The original model became widely known through four key measures: deployment frequency, lead time for changes, change failure rate and time to restore service.


\index{DORA metrics!MTTR}
\index{MTTR}
The current model has five measures and more precise names. DORA explains that the model changed with the technology landscape: \emph{failed deployment recovery time} replaced the broader MTTR label, and \emph{deployment rework rate} became a separate measure. That history matters because dashboards, job interviews and older books still refer to ``the four DORA metrics''. They are describing an earlier version, not an entirely different framework. \autocite{dora-metrics}


DORA organises the current measures under throughput and instability:


\begin{itemize}[leftmargin=*]
\item \textbf{Change lead time} measures the time from a change being committed to version control until it is running in production. Lower is generally better.
\item \textbf{Deployment frequency} measures the number of deployments in a period, or the time between deployments. More frequent delivery can indicate smaller batches and a smoother path to production.
\item \textbf{Failed deployment recovery time} measures how long it takes to recover from a deployment failure that requires immediate intervention. Lower is better.
\item \textbf{Change fail rate} is the proportion of deployments that need immediate intervention, such as a rollback or hotfix. Lower is better.
\item \textbf{Deployment rework rate} is the proportion of unplanned deployments made because of a production incident. Lower is better.
\end{itemize}

The categories are not a management scorecard with one winning number. They describe how changes move and what instability follows. A team deploying daily but performing emergency rework after every third release has produced motion, not dependable delivery.



\subsection{Read Direction And Context, Not A Universal Pass Mark}


Read each metric as a direction of travel for one application or service. A team with changes waiting weeks for a manual test cycle has a different constraint from a team that deploys quickly but spends every afternoon on hotfixes. The first team should investigate the path to production; the second should investigate why planned changes repeatedly create unplanned work.


Do not combine unlike services into a single league table. A mobile consumer application, a safety-critical control system and a quarterly mainframe release operate under different risk, regulation and demand. DORA's current guidance warns that service context matters and that comparisons between unlike applications can mislead. Use the measures to compare a service with its own recent performance and to identify the next bottleneck. \autocite{dora-metrics}


The measurements are also indicators, not proof of causation. DORA reports that they predict organisational performance and team well-being, and its research repeatedly finds that speed and stability are not general trade-offs. Those findings support an improvement hypothesis; they do not prove that increasing a dashboard number will create profit. Delivery performance can improve because teams reduce batch size, automate tests, simplify architecture and learn from failures. Chasing the number while ignoring those capabilities reverses the logic.



\subsection{Why Small Changes Can Improve Speed And Stability}


Return to Meridian's release argument. A monthly release accumulates many changes into one event. When it breaks, the team must search a large set of possible causes, rollback carries more business risk, and the next batch begins queuing while the incident is still being understood.


A team that deploys small changes frequently has a narrower search space. A failed change is easier to identify, revert or repair; automation runs against each small increment; and recovery becomes a routine path rather than an improvised emergency. Frequent deployment is not automatically safe, but small batches, tested rollback and fast feedback can improve both throughput and stability. The pipeline, branching and error-budget practices in the rest of this part show how.


This reframes the release meeting. The useful question is not ``should operations or development win?'' It is ``what evidence shows that this service can move a smaller change safely, and which part of the delivery system currently prevents that?'' The answer may still be to delay a release. The difference is that the delay addresses a named risk rather than preserving a habit.



\subsection{Measure Without Weaponising}


Metrics help only when teams use them as evidence about a delivery system.


Track trends rather than isolated weeks. Correlate delivery data with user outcomes and incident records. Review the five measures together so that apparent throughput cannot hide instability. Pull data from version control, deployment tooling and incident systems where practical, but do not spend months building perfect integrations before holding the first improvement conversation.


Avoid incentives tied to a single measure. Once teams are ranked by deployment frequency, they can split one release into several without improving the service. Once recovery time becomes a target, they can close incidents early. DORA's own guidance warns against setting a metric as the goal, relying on one metric, comparing unlike systems and turning improvement into competition. \autocite{dora-metrics}


The healthier routine is:


\begin{enumerate}[leftmargin=*]
\item Choose one application or service and establish a transparent baseline.
\item Map its delivery path and identify the largest source of waiting or instability.
\item Agree on one improvement, an owner and a review date.
\item Observe all five measures alongside user and operational outcomes.
\item Keep, adapt or reverse the intervention based on the result.
\end{enumerate}

The team owns the interpretation. Managers can remove constraints and fund improvement, but the numbers should not become a device for allocating blame or bonuses.



\subsection{Use The Current Model Precisely}


You should now be able to explain the original four-key model when it appears in older material and use the current five-metric model in new work. More importantly, you should be able to challenge two common errors: treating speed and stability as automatic enemies, and treating a metric as a lever that can be pulled without changing the delivery system underneath it.


For Meridian, the next step is not another argument. It is a service-level baseline, a map of the release path, and a smaller release whose outcome can be observed. Teams that measure delivery performance can improve it deliberately. Teams that reward the appearance of performance teach people to improve the dashboard.

\section{GitHub Actions Workflows}
\index{GitHub Actions}
\index{workflow}
Every team that deploys by hand has a version of the same story. Someone runs the release steps from memory, skips one (clearing a cache, restarting a service, copying one last file) and spends the next four hours working out why production is broken when ``nothing changed.'' The previous topic argued that the cure is a pipeline. This topic is about the most accessible tool for building one: GitHub Actions, the automation platform built into the place your code already lives.


The core idea is simple. You write down your release checklist once, as a file in the repository, and GitHub executes it exactly the same way every time the trigger fires. Fragile human memory becomes reliable machinery. Teams ship faster because nobody is re-typing the same commands, and operations staff see fewer ``works on my machine'' incidents because every change passes through identical automated gates. Just as valuably, every run leaves a record — you always know who deployed what, and when. Automation converts hoping things will work into knowing they will.



\subsection{Anatomy of a workflow}


\index{YAML}
A GitHub Actions \textbf{workflow} is a YAML file stored in your repository under \texttt{.github/workflows/}. YAML is nothing exotic — a human-readable list of instructions that uses indentation for structure, closer to a well-organised shopping list than to a program. Anyone comfortable editing text files can read one, which matters: workflow maintenance is not a developers-only job.


Three concepts do all the work:


\begin{itemize}[leftmargin=*]
\item \textbf{Triggers (events).} A workflow starts when something happens — code is pushed, a pull request is opened, a timer fires on schedule, or a person presses a button.
\item \textbf{Jobs and runners.} Each workflow contains one or more jobs, which execute on \textbf{runners}: temporary virtual machines GitHub provides (or servers you host yourself). Jobs get a fresh machine each run, which is why builds are reproducible.
\item \textbf{Steps.} Inside a job, steps run in order. A step either invokes a prebuilt \textbf{action} (a reusable module published by GitHub, a vendor or the community) or runs a shell command directly.
\end{itemize}

A minimal but genuinely useful workflow looks like this:



\begin{CodeBlock}
name: ci
on: push
jobs:
  test:
    runs-on: ubuntu-latest
    steps:
      - uses: actions/checkout@v4
      - uses: actions/setup-node@v4
        with:
          node-version: 20
      - run: npm ci
      - run: npm test
\end{CodeBlock}


Read it top to bottom: on every push, start an Ubuntu runner, check out the code, install Node.js, install dependencies, run the tests. Eleven lines, and no commit can sneak into the repository without the test suite pronouncing on it.


Because the workflow file lives alongside the code, it goes through the same pull request review as any other change. That is quietly one of the platform's best features: your automation is versioned, visible and improvable. When someone tightens the deployment process, the diff shows exactly what changed and the review shows who agreed to it.



\subsection{What teams actually automate}


Once the basics click, workflows grow to cover the whole delivery path. A common shape: on every push, a matrix job runs the test suite on Windows, macOS and Linux simultaneously — three operating systems tested in the time of one, with nobody lifting a finger. If everything passes, the next job packages the application, uploads the build artifacts, and deploys to a test environment. If any step fails, a notification lands in the team's Slack channel within minutes, so the problem is fixed before customers ever meet it.


The consistency pays a hiring dividend too. A new team member's first contribution goes through the same gates as the tech lead's, with no hidden steps and no tribal knowledge required. The workflow file \emph{is} the documentation of how releases happen — and unlike a wiki page, it can't drift out of date, because it's the thing actually doing the releasing.



\subsection{Deployments, approvals and audits}


\index{AWS}
Deploying to production deserves extra guard-rails, and GitHub provides them through \textbf{environments} — named deployment targets like \texttt{staging} and \texttt{production} that can carry their own protection rules. The most useful rule is a required approval: the workflow runs the build and tests automatically, then pauses at the production gate until a designated person signs off. One click, and the official actions from AWS, Azure or your provider of choice handle the provisioning and release.


\index{change management}
\index{change record}
Every deployment writes a detailed log into the repository's history: which commit, which workflow run, who approved it, when it landed. If you've absorbed Part 2's material on change management, you'll recognise what this is — a change record, generated automatically, that makes audits straightforward instead of painful. A final step can post to Slack or update a ticket so the whole team knows something went live (or didn't) without anyone chasing status. Releases run this way stop being adrenaline events. They become routine — almost boring, which in operations is the highest compliment available.



\subsection{Not just for developers}


It's tempting to file all this under ``programmer concerns,'' but automation reshapes jobs across IT. On a help desk, automated tests and gated deployments mean fewer broken releases landing in your queue at 5 pm Friday. In quality assurance, workflows can spin up disposable test environments on demand, so your hours go into hunting subtle bugs rather than setting up servers. Project managers get a real-time, self-serve view of every release — no more chasing engineers for status. And in a small business, a well-built set of workflows genuinely substitutes for headcount: tasks that once needed a spare pair of hands now run themselves overnight. Whatever role you land in, the pattern holds — automation absorbs the routine so your time goes to judgement, which is the part they actually pay you for.



\subsection{Habits that keep automation healthy}


Workflows accumulate cruft like any other code, so a few disciplines are worth adopting from day one.


\begin{itemize}[leftmargin=*]
\item \textbf{Start small and grow.} Automate the test run first; add packaging, deployment and notifications as confidence builds.
\item \textbf{Keep jobs focused.} One job, one purpose. Small jobs finish faster and make failures easy to pinpoint.
\item \textbf{Apply least privilege.} Grant each job only the access it needs, store credentials in encrypted secrets, and rotate them regularly — a leaked token with narrow permissions is an incident; one with admin rights is a catastrophe.
\item \textbf{Watch run times.} If the five-minute build quietly becomes twenty, feedback is degrading and developers will start skipping it. Treat slowness as a defect.
\item \textbf{Reuse, don't reinvent.} Prefer well-maintained community actions or your own shared modules over bespoke scripts duplicated across repositories.
\item \textbf{Prune periodically.} Review workflows and delete steps that no longer earn their keep.
\end{itemize}


\begin{quote}\small
The least-privilege habit is the one that shows up in post-incident reviews. Workflow logs are public to everyone with repository access, and a secret carelessly echoed into a log is effectively published. Scope tokens tightly, and treat any secret that may have leaked as burned.
\end{quote}



\subsection{The takeaway}


GitHub Actions puts reliable automation within reach of any team — no dedicated build-engineering department required. Because workflows are code, they improve the way code does: incrementally, visibly, with every change reviewed and recorded. Start with something small this week — a workflow that runs your tests on push — and extend it as trust grows. Done well, automation behaves like a colleague who never forgets a step, never gets bored, and works at 3 am without complaint. Give it clear instructions and it will repeat them for you indefinitely — freeing you for the problems that genuinely need a human.

\section{Error Budgets \& SLOs}
\index{error budget}
\index{SLO}
Suppose your company promises customers that its app will be available 99.9\% of the time. That sounds close enough to perfect that most people file it under ``always works.'' Do the arithmetic, though, and 99.9\% allows roughly nine hours of downtime a year — about forty minutes a month. Those forty minutes are not a failure of the promise. They \emph{are} the promise, read from the other side.


\index{SRE}
Site Reliability Engineering (SRE, the discipline Google developed for running planet-scale services) is built on taking that arithmetic seriously. Instead of treating reliability as a vague aspiration (``the site should stay up''), SRE turns it into maths you can track, budget and spend. The two key instruments are service level objectives and the error budgets derived from them, and together they answer the question every product team fights about: when is it safe to keep shipping, and when must we stop and stabilise?



\subsection{Service level objectives}


A \textbf{service level objective (SLO)} is a measurable reliability target for a service: ``99.9\% of requests succeed,'' or ``we respond to user requests within two seconds, 99\% of the time.'' A good SLO has two defining properties. It's expressed as a percentage over a time window, and it measures something \emph{user-facing}. Users don't experience your CPU utilisation; they experience slow pages, failed checkouts and error screens. SLOs are written about the latter, because that's what keeps customers — or loses them.


For non-technical colleagues, the honest translation of an SLO is a promise. Hit it, and users stay happy; miss it, and they notice glitches and slow pages whether or not anyone tells them a number was involved. That's why SLOs are monitored continuously — the point is to notice the service drifting out of the acceptable range before the complaints arrive, not after.



\index{service level agreement}
\begin{quote}\small
Three acronyms travel together and are worth separating once. An \textbf{SLI} (service level indicator) is the measurement itself — the observed success rate this month. An \textbf{SLO} is your internal target for that indicator. An \textbf{SLA} (service level agreement) is the contractual version, the one with refunds attached, which is why SLAs are set looser than SLOs: you want to breach your own target well before you breach the contract.
\end{quote}


The deepest idea in the SLO is choosing a number \emph{below} 100\%. Perfect reliability is unattainable — and long before you got there, each extra ``nine'' would cost more than users would ever notice or pay for. An SLO is an explicit decision about what ``good enough'' looks like: reliable enough to keep customers, honest enough to leave room for change. Because every change (every deployment, migration and config update) carries some risk, a service that tolerates zero failure is a service that can never improve.



\subsection{The error budget}


That gap between perfection and the SLO has a name: the \textbf{error budget}. At 99.9\%, your error budget is the 0.1\% of the window in which the service may fail — those forty-ish minutes a month. And the crucial move is in the word \emph{budget}: this is not shameful slack to be minimised, it's an allowance to be spent.


Every incident, outage and burst of slow responses spends some of it. A twenty-minute outage on Tuesday consumes half the month's budget. What remains tells you, objectively, how much risk you can still afford. Plenty of budget left? Ship the ambitious refactor; run the migration. Budget nearly gone? The data is telling you to slow down. And when the budget is fully spent, a pre-agreed rule kicks in: feature releases pause, and engineering effort pivots to reliability work (fixing the flaky deployment process, adding the missing tests, paying down whatever has been causing the incidents) until the service is back within its SLO.


The elegance of this mechanism is who it aligns. Product managers can see, in hours-remaining terms, how instability eats into feature development time. Engineers know exactly when to slow their release pace, without needing anyone to win an argument. The old standoff (product pushing for speed, operations pleading for caution, decided by whoever shouts loudest) is replaced by a number both sides agreed to in advance.



\subsection{Burn rate: reading the budget over time}
\index{burn rate}


A budget is only useful if you watch the rate at which it's being spent. SRE teams track \textbf{burn rate} — how fast the error budget is being consumed relative to time elapsed in the window. Picture a chart with the weeks of a quarter along the bottom and remaining budget, 100\% down to 0\%, up the side. A healthy service traces a gentle steady slope: routine spending on small incidents, finishing the quarter with budget to spare. A troubled one traces a spiky line — flat stretches punctuated by cliff-drops where a bad release or a cascading failure devours weeks of budget in an afternoon, hitting zero with half the quarter still to run.


Teams commonly mark zones on that chart and attach pre-agreed responses:


\begin{itemize}[leftmargin=*]
\item \textbf{Green (healthy):} plenty of budget for the time elapsed — keep shipping normally.
\item \textbf{Amber (caution):} burning faster than the window allows — slow down, defer the risky changes, add extra review and testing.
\item \textbf{Red (freeze):} budget exhausted or nearly so — stop feature releases and put engineering effort into recovery and prevention.
\end{itemize}

The point of pre-agreeing the zones is that nobody has to improvise policy during a bad week; the chart makes the call. A fast burn is also diagnostic. It means one of two things: your releases are riskier than your process can absorb, or your SLO is stricter than your architecture can honestly support. Both are fixable — tighten testing and slow deployments in the first case; renegotiate the SLO in the second, if the customer impact genuinely warrants it. An SLO is a decision, and decisions can be revisited — openly, with data, rather than by quietly ignoring the target.



\subsection{What error budgets change}


Notice what has happened to the conversation. ``Is the service reliable enough?'' used to be answered with anecdotes and anxiety. With SLOs and error budgets it's answered with a number, and the follow-up — ``so can we ship this week?'' — answers itself. The framework takes emotion out of decisions that used to be political, which is why it has spread far beyond Google.


For you, entering the industry, the practical takeaway is twofold. First, reliability is a feature with a cost and a target (not an absolute) and you should be suspicious of any team that claims to aim for 100\%, because they're really aiming for ``we haven't decided.'' Second, when you're under budget you have earned the freedom to move fast, and when the budget is nearly spent, stability \emph{is} the roadmap. Teams that internalise this stop having the speed-versus-reliability argument entirely. They had it once, wrote the answer down as an SLO, and let the budget referee every week since.

\section{Trunk-Based Development vs Feature Branching}
\index{feature branching}
\index{trunk-based development}
\index{DORA metrics}
At Meridian Insurance (the firm from the DORA metrics topic) a contractor once spent six weeks building a quoting feature on his own branch. The work was good. The merge was not. In the six weeks his branch had been drifting away from the main line, four other teams had reshaped the modules he depended on. Reconciling the two histories took nine days, broke two things the tests didn't cover, and produced the sentence every engineering manager dreads: ``it worked before the merge.''


How a team uses branches sounds like a low-stakes tooling preference. It isn't. Branching strategy determines how quickly changes integrate with everyone else's work, and therefore how much merge pain you accumulate — which makes it one of the strongest levers on the DORA metrics you've just met. There are two poles to the debate, and most real teams live somewhere between them.



\subsection{Trunk-based development}


\index{CI/CD}
In trunk-based development, everyone commits to a single shared main line (the trunk) in small, frequent increments, at least daily. There are no long-lived personal branches accumulating private history; integration happens continuously, which is what ``continuous integration'' originally meant before it became the name of a build server.


\index{feature flags}
The immediate objection is obvious: what about half-finished work? You can't ship half a quoting feature to production. The trunk-based answer is the \textbf{feature toggle} (or feature flag): the incomplete code merges into trunk and deploys with everything else, but sits behind a switch that keeps it invisible to users until it's ready. The code integrates early; the \emph{feature} launches when you choose. Toggles bring their own housekeeping (old flags need removing, and testing must cover both switch positions) but they decouple two things that were never naturally coupled: integrating code and releasing features.


What trunk-based development buys you is the near-elimination of merge conflicts. When everyone integrates daily, nobody's view of the codebase is more than a day stale, so collisions are small and caught while both authors still remember what they wrote. What it demands is discipline: a solid automated test suite (because trunk must stay releasable), the skill of slicing features into small safe increments, and a team that treats a broken build as everyone's top priority.



\subsection{Feature branching}


Feature branching is the strategy many newcomers meet first, because it is the default shape of working with GitHub. Each piece of work gets its own isolated branch; when it is finished, a pull request proposes the merge back, a colleague reviews the diff, the automated checks run, and the branch lands.


\index{open source}
Isolation is genuinely valuable. You can experiment without destabilising anyone, park work when priorities shift, and (most importantly) the pull request gives code review a natural home. The review conversation attached to a PR is often the best design documentation a change ever gets, and for open source projects, where contributors are strangers, this model is essentially mandatory.


The cost is drift. Every day a branch lives apart from trunk, two things grow silently: the branch's ignorance of what has changed underneath it, and trunk's ignorance of what the branch is about to do to it. Merge conflicts are the visible symptom; the invisible one is worse — two branches that don't textually conflict but break each other's assumptions, sailing through the merge and failing in production. Integration pain isn't eliminated by branching; it's deferred, and it compounds while it waits. Meridian's nine-day merge wasn't bad luck. It was six weeks of deferred integration, paid back with interest.



\subsection{Choosing — and blending}


Framed as a binary, the choice looks hard: trunk-based development maximises flow and feedback; feature branches give you isolation and reviewable units of work. In practice, almost every effective team blends the two, and the blend is less a compromise than a synthesis.


The recipe: keep branches, but keep them \emph{short-lived}. A branch carries one small piece of work, lives a day or two, and merges via a pull request that a colleague can actually review in ten minutes — because it's a two-hundred-line diff, not a six-week epic. Larger features are sliced into a sequence of such branches, integrated one at a time behind a toggle if needed. You get the review culture and safety of pull requests with integration frequency close enough to trunk-based to keep merges trivial.


The evidence backs the blend. The \emph{Accelerate} research found that teams with branches lasting less than about a day, and only a handful active at once, showed higher delivery performance — while long-lived branches were associated with slower, less stable delivery. The finding, notably, is about branch \emph{lifetime}, not about whether branches exist at all.



\begin{quote}\small
A useful team norm: if a branch is about to see its third sunrise, something is wrong — slice the work smaller, or merge what's safe behind a toggle. Branch age is one of the cheapest early-warning metrics a team can watch.
\end{quote}


It's the mirror image of the DORA logic from earlier in this part: many small deployments beat one big deployment, and many small merges beat one big merge, for the same reason. Small changes are easy to understand, easy to review, easy to test and easy to reverse.



\subsection{The takeaway}


Don't get religious about the label; get disciplined about the interval. Whether your team calls its process trunk-based or feature branching matters far less than how long changes stay isolated from the main line. Keep merges small and frequent, keep branches measured in hours or days rather than weeks, use toggles to separate integration from release, and reserve long-lived branches for the rare cases that truly need them. When you join a team, you can read its health in its repository: dozens of stale branches and dreaded ``merge parties'' mean deferred integration debt; a busy trunk fed by short-lived branches means the debt is being paid daily, in small instalments, while it's still cheap.

\section{Practice Artefact}

\index{change failure rate}
\index{CI/CD}
\index{DORA metrics!change failure rate}
\index{DORA metrics!deployment frequency}
\index{DORA metrics!lead time for changes}
Produce a pipeline evidence sheet. Build or sketch a three-step CI/CD pipeline, record deployment frequency, lead time for changes, change failure rate, and time to restore for one deliberate failure, then write the one-paragraph improvement you would try next.


\index{workflow}
The artefact should include the workflow file or pseudocode, the failed run, the recovery step, and a short explanation of why the change was safe enough to automate.
